\section{Donor-engulfment architecture}

V3 explicitly imports the strongest useful structure from two neighboring layers.

\subsection{Action, delegation, and authorization donors}

P8 reuses evidence-backed permission graphs and external deterministic authorizers; runtime-neutral action certificates carrying action identity, approvals, runtime and outcome receipts; append-only human-to-agent delegation provenance; authority-narrowing and cascade-containment rules; principal-chain composition with bounded scope; typed verifier certificates; heterogeneous authorization-evidence chains; and generic cross-domain authority relations.

These mechanisms answer important local questions such as whether an action is permitted, whether a delegation is valid, whether required evidence certificates are present, or whether heterogeneous authorization receipts jointly satisfy an action policy. This paper claims none of those mechanisms as its own.

\subsection{Scientific claim, release, and adjudication donors}

P8 also reuses domain scientific harnesses that separate semantic from scientific authority and authorize evidence-bound release; claim-to-evidence chains with deterministic grounding and verification; persistent contract-governed research artifacts; evidence-calibrated claim adjudication; and evidence-ledger review that routes unsupported, contradicted, or mixed-evidence claims back for revision.

These donors establish that scientific claim validation/release can itself be externalized and governed. P8 therefore does not claim novelty merely for separating a model's fluent answer from a scientific release decision.

\subsection{The remaining composition problem}

The donor mechanisms are powerful but heterogeneous. Their local verdicts concern different objects: an action, a delegation, a receipt chain, a verifier predicate, a claim-evidence packet, a research artifact, or a domain-specific release decision. This paper's V3 object is the common relation that governs when such a local authority object is entitled to discharge a target scientific obligation of another type or scope.
